\section{INFRAESTRUTURA}
\label{INFRA}
Para execução das atividades do \NOMECURSO está contemplada a seguinte infraestrutura:

\subsection{Laboratórios}

\begin{itemize}%[LAMEP]
	\item [LMAT] LABORATÓRIO DE MATERIAIS: São realizadas análises microestruturais, tratamentos térmicos, ensaios mecânicos e ensaios de
	corrosão. Disciplinas contempladas: \LMATdisc.
	
	\item [LIAF] LABORATÓRIO DE INSPEÇÃO E ANÁLISE DE FALHAS: São realizados ensaios não-destrutivos com equipamentos tais como
	ultra-som, kit de inspeção por partículas magnéticas, kit de inspeção por líquido penetrante, endoscópio industrial,
	viscosímetro e espectrômetro de fluorescência de raios-X. Disciplinas contempladas: \LIAFdisc.
	
	\item [LMET] LABORATÓRIO DE METROLOGIA DIMENSIONAL: São realizadas medições com uso de paquímetros digitais e analógicos, medições
	com uso de micrômetros digitais e analógicos, medições com uso de goniômetros, medições com uso de relógios
	comparadores digitais e analógicos, medições com uso de rugosímetros, medições com projetor de perfil, medições com uso
	de traçadores de altura, medições com uso de réguas graduadas, medições e engenharia reversa com tecnologia ótica 2D e
	medições por coordenadas e engenharia reversa com tecnologia de scanner a laser.  Disciplinas contempladas: \LMETdisc.
	
	\item [LTF] LABORATÓRIO DE MÁQUINAS TÉRMICAS E DE FLUXO: Laboratório cuja finalidade é propiciar aos alunos uma visualização prática
	do princípio de funcionamento das máquinas de fluxo e térmicas. Também ensaiam-se experimentos práticos e
	computacionais que servirão de base para produção científica e desenvolvimento de mão de obra especializada. Disciplinas contempladas:\LTFdisc.
	
	\item [LIA] LABORATÓRIO DE INFORMÁTICA APLICADA: Tem por finalidade desenvolver atividades voltadas à simulação de circuitos
	elétricos e eletrônicos, desenvolvimento de algoritmos computacionais, bem como a criação e desenvolvimento de modelos
	em 2D e 3D. Disciplinas contempladas: \LIAdisc.
	
	\item [LSHIP] LABORATÓRIO DE SISTEMAS HIDRÁULICOS E PNEUMÁTICOS: Tem por finalidade possibilitar que os alunos conheçam os principais
	atuadores pneumáticos e hidráulicos, dispositivos de geração e tratamento de ar comprimido, simular e implementar de
	forma experimental circuitos pneumáticos e hidráulicos. Disciplinas contempladas: \LSHIPdisc.
	
	\item [LAMEP] LABORATÓRIO DE ACIONAMENTOS DE MÁQUINAS E ELETRÔNICA DE POTÊNCIA: Este laboratório dispõe de equipamentos e dispositivos
	que permitem a realização de aulas práticas em circuitos elétricos em corrente alternada (cargas R, RL, RC, RLC),
	comandos e instalações elétricas em âmbito industrial, acionamentos e controladores industriais aplicados em motores de
	indução monofásicos, trifásicos e servomotores CA, com e sem carga. Disciplinas contempladas:\LAMEPdisc.
	
	\item [LEE] LABORATÓRIO DE ELETROELETRÔNICA: Este laboratório dispõe de equipamentos e dispositivos que permitem a realização de aulas práticas em circuitos elétricos em corrente contínua, bem como medições elétricas das principais grandezas (resistência, tensão, corrente, etc), módulos para práticas em circuitos de eletrônica digital e analógica, envolvendo amplificadores operacionais, diodos, transistores, etc. Disciplinas contempladas:\LEEdisc.
	
	\item [LINC] LABORATÓRIO DE INSTRUMENTAÇÃO E CONTROLE: Tem por finalidade a realização de atividades e projetos da área de instrumentação e controle, através da utilização de equipamentos que simulam situações reais do ambiente industrial. Através~ de modernos sistemas eletrônicos de calibração e medição de variáveis de processo como nível, pressão, vazão e temperatura. Possibilita ao aluno associar a teoria aos conhecimentos práticos e reais da instrumentação industrial.
	Disciplinas contempladas:\LINCdisc.
	
	%\item [LPROT] LABORATÓRIO DE PROTÓTIPOS: Tem por finalidade permitir ao aluno a realização de projetos de protótipos de processos elétricos e eletrônicos, através da simulação~ em~ escala real ou reduzida, possibilitando uma posterior utilização em ambientes industriais ou sistemas embarcados.
	%    Disciplinas contempladas:\LPROTdisc.
	
	\item [LPC] LABORATÓRIO DE POTÊNCIA E CONTROLE: Tem por finalidade permitir ao aluno a realização de projetos de protótipos de processos elétricos e eletrônicos, através da simulação e concepção em escala real ou reduzida, possibilitando uma posterior utilização em ambientes industriais ou sistemas embarcados.
	Disciplinas contempladas:\LPCdisc.
	
	\item [LAMSC] LABORATÓRIO DE APRENDIZAGEM DE MÁQUINAS E SIMULAÇÃO COMPUTACIONAL: Tem o objetivo de permitir ao aluno implementar e testar modelos computacionais para aplicações de aprendizagem de máquinas em problemas de controle e identificação de sistemas. Além disso, também possibilita ao aluno a montagem de projetos de protótipos em microcontroladores e robótica.
	Disciplinas contempladas:\LAMSCdisc.
	
	
	%\item [LEE] LABORATÓRIO DE ELETROELETRÔNICA: Este laboratório dispõe de equipamentos e dispositivos que permitem a realização de
	%aulas práticas em circuitos elétricos em corrente contínua, bem como medições elétricas das principais grandezas
	%(resistência, tensão, corrente, etc), módulos para práticas em circuitos de eletrônica digital e analógica, envolvendo
	%amplificadores operacionais, diodos, transistores, etc. Disciplinas contempladas: Eletricidade I, Eletrônica Analógica
	%e Eletrônica Digital.
	%
	%\item [LINC] LABORATÓRIO DE INSTRUMENTAÇÃO E CONTROLE: Tem por finalidade a realização de atividades e projetos da área de
	%instrumentação e controle, através da utilização de equipamentos que simulam situações reais do ambiente industrial.
	%Através~ de modernos sistemas eletrônicos de calibração e medição de variáveis de processo como nível, pressão, vazão e
	%temperatura, torna o aluno capaz de associar a teoria~ aos conhecimentos práticos e reais da instrumentação industrial.
	%Disciplinas contempladas: Instrumentação e Controle, Redes Industriais e CLP.
	%
	%\item [LPROT] LABORATÓRIO DE PROTÓTIPOS: Tem por finalidade permitir ao aluno a realização de projetos de protótipos de processos
	%elétricos e eletrônicos, através da simulação~ em~ escala real ou reduzida, possibilitando uma posterior utilização em
	%ambientes industriais ou sistemas embarcados. Contempla todas as disciplinas do curso. (J.DANIEL)
	
\end{itemize}

\subsection{Ambientes Administrativos}

%A infraestrutura de salas de aula e laboratórios, %, além de complementar a teoria aprendida em sala de aula e laboratórios, também irá auxiliar
Os alunos do \NOMECURSO contam com uma infraestrutura administrativa composta de: % no desenvolvimento de atividades de pesquisa e extensão.

\begin{itemize}
	\item SALA DE PROFESSORES
	\item SALA DE ATENDIMENTO AO ALUNO
	\item SALA DA COORDENADORIA DOS CURSOS DO EIXO TECNOLÓGICO DA INDÚSTRIA
	\item 5 SALAS DE AULA
	\item GABINETE DE PROFESSORES
	\item AUDITÓRIO
	\item SALA DE VIDEOCONFERÊNCIA
	\item BIBLIOTECA
	\item CENTRO ACADÊMICO DO CURSO
\end{itemize}

Essa infraestrutura também auxilia os alunos do \NOMECURSO no desenvolvimento de atividades de pesquisa e extensão.

\subsection{Bloco II do Eixo Tecnológico da Indústria}
O Bloco II do Eixo Tecnológico da Indústria, denominado Centro de Pesquisa e Tecnologia (CPT), conta com os seguintes ambientes:
\begin{itemize}
	\item \textbf{Finalizados:} ajustagem, soldagem, manutenção eletromecânica, usinagem convencional, usinagem CNC, Laboratório de CAD e CAM;
	\item \textbf{Em construção:} 08 (oito) laboratórios de pesquisa e desenvolvimento, 15 (quinze) gabinetes para professores, biblioteca setorial, auditório e salas de aula para futuros
	programas de pós-graduação.
\end{itemize}

Essa nova infraestrutura contribuirá de maneira extremamente significativa para a constante
atualização tecnológica do \NOMECURSO. Maiores detalhes podem ser apreciados no projeto do CPT apresentado no ANEXO \ref{anexoProjeto}.
